\section{Two posterior-specific reductions}
We first quantify how the posterior relates to convex path-norm
regression on the observed sample.
Let
\[
 z(x)=\sum_{j=1}^N v_j(a_jx+b_j)_+,
 \quad (a_j,b_j)\sim N(0,\nu I_2),\quad
 v_j\sim N(0,\nu/N^2),\quad \omega=N/\nu,
\]
where all coordinates are independent and $0<\nu\le1$.
The actual posterior has density proportional to
$\exp[-\omega\widehat R(z)/(2\lambda)]$ relative to this prior.
Throughout this section $|X_i|,|Y_i|\le1$.

Write $(a_j,b_j)=r_j\theta_j$ with $\theta_j$ uniform on the unit
circle, and put $c_j=v_jr_j$. The exact pushforward prior has independent
uniform directions and independent Laplace amplitudes
$c_j\sim\operatorname{Laplace}(0,\omega^{-1})$.
Indeed $r_j^2$ is exponential of mean $2\nu$, and
\[
 \mathbb E e^{itc_j}
 =\mathbb E\exp[-\nu t^2r_j^2/(2N^2)]
 =(1+\nu^2t^2/N^2)^{-1}.
\]
Put $P(c)=\sum_j|c_j|$ and
$A_\theta c=(\sum_jc_j(\theta_{j,1}X_i+\theta_{j,2})_+)_{i=1}^n$.
All statements below are about this exact pushforward posterior;
no replacement prior or estimator is introduced.

\begin{proposition}[Expected excess mass on a training-prediction fiber]
Fix the directions and a feasible training vector $b$.
Let $d_\theta=N-\operatorname{rank}(A_\theta)$ and
\[
 P_{\min}(b,\theta)=\min\{\|c\|_1:A_\theta c=b\}.
\]
For almost every fiber under the regular conditional posterior,
\[
 \mathbb E[P(c)-P_{\min}(b,\theta)\mid\theta,A_\theta c=b]
 \le d_\theta/\omega\le\nu. \tag{PR1}
\]
\end{proposition}
\begin{proof}
Conditional on the fiber, the likelihood is constant. The posterior
density with respect to its affine Hausdorff measure is proportional
to $e^{-\omega\|c\|_1}$; the coarea normalization is constant there.
Choose a minimizer $c_*$. In orthonormal nullspace coordinates $u$,
let $V(u)=\|c_*+Bu\|_1$, where $B$ is an isometry from
$\mathbb R^{d_\theta}$ onto $\ker A_\theta$.
The density is coercive and has exponential tails. Integration of
$\operatorname{div}(u e^{-\omega V(u)})$ gives
\[
 d_\theta=\omega\mathbb E[u\cdot\nabla V(u)].
\]
The piecewise-linear convex potential is differentiable almost
everywhere; the identity also follows by smooth approximation.
Convexity gives $u\cdot\nabla V(u)\ge V(u)-V(0)$.
The zero-dimensional case is immediate. This proves (PR1).
\end{proof}

\begin{proposition}[Approximate minimization of the full empirical path objective]
Suppose $N\ge n$ and $0<\lambda\le1/2$. Let $J_*$ be the minimum,
over all finite shallow ReLU implementations, of
\[
 J(c,\theta)=\widehat R(z_{c,\theta})+2\lambda\sum_j|c_j|.
\]
Equivalently, one may first minimize the amplitude mass among
implementations of each training vector, then minimize over those vectors. Then the actual posterior satisfies
\[
 0\le\mathbb E_\Pi[\widehat R(z)+2\lambda P(c)]-J_*
 \le C\lambda\nu\log\frac{e}{\lambda\nu}. \tag{PR2}
\]
Here $C$ is universal. The mass in the posterior expectation is that
of its actual amplitude representation, not a separately minimized
version.
\end{proposition}
\begin{proof}
The signed unit-ridge training vectors form a compact subset of
$\mathbb R^n$. The minimum-mass representation of any feasible
training vector therefore exists, and a linear-dependence reduction
uses at most $n$ atoms. Minimizing the coercive objective over the
resulting compact sublevel sets gives a minimizer $f_*$ with at most
$n$ atoms. Since zero is feasible,
\[
 J_*\le1,\qquad \widehat R(f_*)\le1,\qquad P_*\le(2\lambda)^{-1}.
\]
Thus $f_*$ fits into the prescribed width $N$.

Let $F=-(2\lambda/\omega)\log Z$ be the actual free energy.
Represent $f_*$ by the balanced Gaussian center
$\bar r_j=\sqrt{N|c_j|}\theta_j$,
$\bar v_j=\operatorname{sign}(c_j)\sqrt{|c_j|/N}$.
Its weighted Gaussian energy is $2P_*$.
Around this center use native covariance multiplied by
$\epsilon^2$, with $\epsilon=(\lambda\nu)^2$.
The comparison law has
\[
 \operatorname{KL}(\rho\Vert\mu)
 \le\omega P_*+3N\log(1/\epsilon).
\]
Its uniform root-mean-square function perturbation is bounded by
$C\epsilon\sqrt{\nu P_*}+C\epsilon^2\nu$.
Using $\widehat R(f_*)\le1$ gives
$\rho\widehat R\le\widehat R(f_*)+C\lambda\nu$.
The Gibbs variational inequality therefore yields
\[
 F\le J_*+C\lambda\nu\log\frac e{\lambda\nu}. \tag{PR3}
\]

For the pushforward prior,
$\mathbb E e^{t\omega P}=(1-t)^{-N}$ for $0<t<1$.
The entropy inequality and the Gibbs identity imply
\[
 F\ge\mathbb E_\Pi\widehat R+2\lambda t\mathbb E_\Pi P
       +2\lambda\nu\log(1-t).
\]
Radial integration by parts in the Laplace amplitudes gives
\[
 \mathbb E_\Pi P
 =\nu-\lambda^{-1}\mathbb E_\Pi P_n[z(z-Y)]
 \le\nu+\widehat{Y^2}/(4\lambda)
 \le\nu+(4\lambda)^{-1}.
\]
Use $1-t=\lambda\nu$. The lost coercivity term
$2\lambda(1-t)\mathbb E_\Pi P$ is at most $C\lambda\nu$.
Combining this with (PR3) proves the upper bound in (PR2).
Its lower bound follows from the definition of $J_*$.
\end{proof}

Neither statement alone controls population predictions in an empty
sample interval. In (PR1), a small mass excess need not imply small
function distance to a minimum-mass representative: nearby ridge
directions can make the relevant stability constant arbitrarily bad.
In (PR2), near-minimal objective value is not a uniform output-range
bound. The four-point example later in this paper proves that
full Gaussian path-norm interpolants themselves can have arbitrarily
large excursions between bounded training labels. A statistical
improvement must use the local sample spacings and the cost of fitting
nearby conflicting labels, rather than invoking deterministic range
preservation.
